\section{KKT 把最优性写成一组方程}
\label{sec:09-Convex-Optimization/05-Duality/03}

你有一个候选点。梯度看着挺小，约束也差不多满足。但``差不多''不是证明。你需要的是一组方程——每一项能逐一核对，全部通过则最优性成立，任意一条失败就告诉你改进方向。

当强对偶成立且原始、对偶最优点都能取得时，Lagrangian 推导链中的每一步不等式在最优点必须全部取等号。把这些等号条件逐条摘出来，就是 Karush--Kuhn--Tucker（KKT）条件。

\subsection{四组条件，各自把关一个环节}

对可微问题

\[
\begin{aligned}
\text{minimize}\quad&f_0(x)\\
\text{subject to}\quad&f_i(x)\le0,\quad i=1,\ldots,m,\\
&h_j(x)=0,\quad j=1,\ldots,p,
\end{aligned}
\]

KKT 条件包括四组：

\textbf{第一组：原始可行性}

\[
f_i(x^\star)\le0,\qquad
h_j(x^\star)=0.
\]

候选点必须落在可行域内。这是最基本的入场券——不等式不能违反，等式必须精确成立。任何一条不满足，这个点甚至不配称为候选解。

\textbf{第二组：对偶可行性}

\[
\lambda_i^\star\ge0.
\]

不等式的 Lagrange 乘子不能为负。原因是 Lagrangian

\[
L(x,\lambda,\nu)=f_0(x)+\sum_i\lambda_i f_i(x)+\sum_j\nu_j h_j(x)
\]

中，乘子 \(\lambda_i\) 乘以约束值 \(f_i(x)\)。若 \(\lambda_i<0\)，那么当约束被违反（\(f_i(x)>0\)）时，\(\lambda_i f_i(x)\) 是一个负数——惩罚变成了奖励——Lagrangian 就不再是原目标的一个有效下界，整个对偶框架的上确界构造就会失效。

\textbf{第三组：互补松弛}

\[
\lambda_i^\star f_i(x^\star)=0
\qquad\text{对每个 }i.
\]

对每一个约束 \(i\)，\(\lambda_i^\star\) 和 \(f_i(x^\star)\) 至少有一个是零，不允许两者同时为正。这条条件来自强对偶的等号链条。在 \(x^\star,\lambda^\star,\nu^\star\) 处：

\[
f_0(x^\star)=d^\star=g(\lambda^\star,\nu^\star)\le L(x^\star,\lambda^\star,\nu^\star)=f_0(x^\star)+\sum_i\lambda_i^\star f_i(x^\star).
\]

但 \(\lambda_i^\star\ge0\)，\(f_i(x^\star)\le0\)，所以 \(\sum_i\lambda_i^\star f_i(x^\star)\le0\)。要使上面的不等号成为等号，每一项 \(\lambda_i^\star f_i(x^\star)\) 都必须等于零。

\textbf{第四组：驻点条件}

\[
\nabla f_0(x^\star)
+\sum_i\lambda_i^\star\nabla f_i(x^\star)
+\sum_j\nu_j^\star\nabla h_j(x^\star)
=0.
\]

这是 Lagrangian \(L(\cdot,\lambda^\star,\nu^\star)\) 关于 \(x\) 的梯度在 \(x^\star\) 处为零。它的几何含义：在可行域内部（满足等式约束的切空间内），目标梯度被活跃约束的梯度合力抵消，任何可行方向上的移动都无法改进目标值。换句话说，目标函数的下降方向全部被约束的梯度方向覆盖了——你无法在满足约束的前提下挪到更优的点。

\subsection{互补松弛在区分两种约束}

互补松弛的乘积为零这一形式，将约束分为两类。

如果 \(f_i(x^\star)<0\)——约束在最优点严格满足，还有余量——互补松弛迫使 \(\lambda_i^\star=0\)。这个约束在当前位置是松弛的，它的存在对一阶最优性没有影响，其边际价格为零。想象一个约束要求 \(x\le 10\)，而最优解正好在 \(x=3\)。约束``\(x\le 10\)''在 \(x=3\) 处是不活跃的——你不会为了再多一单位的 \(x\) 而愿意付出任何代价，因为现在已经远在边界以内。

如果 \(\lambda_i^\star>0\)，则必须 \(f_i(x^\star)=0\)。约束紧贴边界，乘子的大小衡量了``放松约束一个单位能带来多少目标改善''。\(\lambda_i^\star\) 在经济学中就是约束 \(f_i(x)\le0\) 的影子价格。

但反向不成立。一个约束可以 \(f_i(x^\star)=0\) 同时 \(\lambda_i^\star=0\)。比如两个约束写的是同一个条件（重复约束），或者该约束虽然在边界上却与目标函数的下降方向恰好正交——它在当前几何结构中刚好不是限制性的。``活跃''（取到边界）和``有正价格''（乘子大于零）是两种相关但不等价的性质。互补松弛说的是乘积为零，不要机械地在二者之间做双向推断。

\subsection{一维例子把四组条件走一遍}

\[
\begin{aligned}
\text{minimize}\quad&x^2\\
\text{subject to}\quad&1-x\le0.
\end{aligned}
\]

写下 KKT：

\begin{itemize}
\tightlist
\item
  原始可行性：\(1-x\le0\)，即 \(x\ge1\)。
\item
  对偶可行性：\(\lambda\ge0\)。
\item
  互补松弛：\(\lambda(1-x)=0\)。
\item
  驻点条件：\(2x-\lambda=0\)。
\end{itemize}

试探约束不活跃：若 \(1-x<0\)（即 \(x>1\)），互补松弛迫使 \(\lambda=0\)，驻点条件得 \(x=0\)。但 \(x=0\) 不满足 \(x\ge1\)——原始可行性失败。因此约束必须活跃：\(x^\star=1\)。代入驻点条件：\(\lambda^\star=2\)。

这个解有两层对应关系，而且两层对应都能在一张图上看到。

几何上：无约束目标 \(x^2\) 的最低点在 \(x=0\)，梯度指向 \(x\) 增大的方向（在 \(x=1\) 处梯度为 \(+2\)，指向 0 的方向）；边界 \(x\ge1\) 把 \(x\) 挡在 1 处不让你再往左走。乘子 \(\lambda=2\) 给出了一个法向力——大小 2，方向指向可行域内部（\(x\) 增大的方向），恰好抵消目标梯度。目标把你向左推，约束把你向右推，合力为零——这就是驻点条件的几何。

对偶上：Lagrangian \(L(x,\lambda)=x^2+\lambda(1-x)=\lambda + (x^2-\lambda x)\)。在最优乘子 \(\lambda^\star=2\) 下，\(L(x,2)=2+x^2-2x=(x-1)^2+1\)。这个函数在 \(x=1\) 处取最小值 1，恰好等于 \(f_0(1)=1\)。对偶间隙闭合。

再做一个带等式约束的例子。

\[
\begin{aligned}
\text{minimize}\quad&x_1^2+x_2^2\\
\text{subject to}\quad&x_1+x_2=1.
\end{aligned}
\]

这是从原点到直线 \(x_1+x_2=1\) 的最短距离。KKT 条件：

\begin{itemize}
\tightlist
\item
  原始可行性：\(x_1+x_2=1\)。
\item
  驻点条件：\(\nabla(x_1^2+x_2^2)+\nu\nabla(x_1+x_2-1)=0\)，即 \(2x_1+\nu=0\)，\(2x_2+\nu=0\)。
\end{itemize}

驻点给出 \(x_1=x_2=-\nu/2\)。代入等式约束：\(x_1+x_2=-\nu=1\)，所以 \(\nu=-1\)，\(x_1=x_2=0.5\)。最优点在直线中点，验证了点到直线垂足的几何事实。注意这里没有不等式约束，互补松弛和对偶可行性不参与（或者说对等式乘子 \(\nu\) 没有符号约束）。

\subsection{凸问题中 KKT 是充要条件}

若 \(f_0,f_i\) 凸且等式约束为仿射（\(h_j(x)=a_j^{\trans}x-b_j\)），则任何满足 KKT 的点都是全局最优。充分性的证明只需把 Lagrangian 下界链在最优乘子处收紧。

固定 \(\lambda^\star,\nu^\star\)。由于 \(\lambda_i^\star\ge0\) 且各 \(f_i\) 凸，仿射项保持凸性：

\[
L(x,\lambda^\star,\nu^\star)=f_0(x)+\sum_i\lambda_i^\star f_i(x)+\sum_j\nu_j^\star(a_j^{\trans}x-b_j)
\]

是 \(x\) 的凸函数（凸函数的非负组合加仿射函数仍为凸）。驻点条件说它在 \(x^\star\) 处梯度为零：

\[
\nabla_x L(x^\star,\lambda^\star,\nu^\star)
=\nabla f_0(x^\star)
+\sum_i\lambda_i^\star\nabla f_i(x^\star)
+A^{\trans}\nu^\star=0.
\]

可微凸函数的一阶条件：若梯度在某点为零，则该点是全局最小值点。因此：

\[
L(x^\star,\lambda^\star,\nu^\star)\le L(y,\lambda^\star,\nu^\star)
\quad\text{对所有 }y.
\]

现在把链条串起来。对任意原始可行 \(y\)（即 \(f_i(y)\le0\) 且 \(h_j(y)=0\)）：

\[
\begin{aligned}
f_0(x^\star)
&= f_0(x^\star)+\sum_i\lambda_i^\star f_i(x^\star)+\sum_j\nu_j^\star h_j(x^\star)
\quad\text{(互补松弛 + 等式约束)}\\
&= L(x^\star,\lambda^\star,\nu^\star)\\
&\le L(y,\lambda^\star,\nu^\star)
\quad\text{(凸函数全局最小)}\\
&= f_0(y)+\sum_i\lambda_i^\star f_i(y)+\sum_j\nu_j^\star h_j(y)\\
&\le f_0(y)
\quad\text{(\(\lambda_i^\star\ge0,\;f_i(y)\le0,\;h_j(y)=0\))}.
\end{aligned}
\]

每一步不等号的依据都直接来自 KKT 的某一条。没有哪一步是靠``猜测''或``近似''——这是一条全等号的证明链。

反方向上，若凸问题还满足 Slater 条件（或更一般的约束资格），KKT 也是最优性的必要条件——任何最优点必定存在一组乘子满足 KKT。因此在``凸 + Slater''的前提下，求解原问题等价于寻找满足 KKT 四组方程的 \((x^\star,\lambda^\star,\nu^\star)\)，这是一个从优化问题到方程组求解的降维。

\subsection{离开凸性后 KKT 大幅缩水}

非凸问题中，KKT 只能用来筛选候选点。它不再是充分条件，也不再总是必要条件。

驻点条件 \(\nabla L=0\) 对一般的非凸函数只保证``梯度方向无法再改进''——它可能是局部最小、局部最大、鞍点，或某个曲率更复杂的临界点。考虑一个二维非凸例子：

\[
\begin{aligned}
\text{minimize}\quad&x_1^2-x_2^2\\
\text{subject to}\quad&-1\le x_1\le 1,\;-1\le x_2\le 1.
\end{aligned}
\]

目标函数在原点 \((0,0)\) 处满足驻点条件——梯度 \(2x_1,-2x_2\) 为零。但在 \(x_2\) 方向它是凹的，沿着 \(x_2=\pm 1\) 移动目标可以无限下降（在框约束内下到 \(-1\)）。原点是一个鞍点，满足 KKT 却绝对不是最优——连局部最优都不是。

更隐蔽的问题是约束资格失败。假设在最优点处，两个不等式约束的梯度恰好线性相关。此时可能不存在任何实数乘子对 \((\lambda_1,\lambda_2)\) 同时满足驻点条件。求解器返回``未能收敛''，但根因不在算法——而在于这个问题在那个点上本身就是 KKT 条件退化的。

对于不可微凸问题，驻点条件需要写成次微分形式：

\[
0\in
\partial f_0(x^\star)
+\sum_i\lambda_i^\star\partial f_i(x^\star)
+A^{\trans}\nu^\star,
\]

其中 \(\partial f\) 是凸函数 \(f\) 的次微分集合（所有次梯度的集合）。但次微分的加法需要正则条件——至少需要一个函数在相关点的相对内部连续——否则 \(\partial(f+g)\) 可能只是 \(\partial f+\partial g\) 的真子集。使用这条公式前应查阅凸分析中次微分 calculus 的精确表述，而非随手写成。

\subsection{KKT 残差是数值求解后的体检表}

求解器返回候选 \((\hat x,\hat\lambda,\hat\nu)\) 后，逐项核算以下残差：

\begin{itemize}
\tightlist
\item
  原始不等式残差：\(\max(f_i(\hat x),0)\)——每个不等式的违反量，只看正部，因为\(f_i\le0\)才算满足；
\item
  原始等式残差：\(|h_j(\hat x)|\)——等式偏离零的程度；
\item
  对偶残差：\(\max(-\hat\lambda_i,0)\)——乘子为负的大小，负值是不允许的；
\item
  互补残差：\(|\hat\lambda_i f_i(\hat x)|\)——每个约束的乘积的绝对值，两者应该至少有一个接近于零；
\item
  驻点残差：\(\|\nabla_x L(\hat x,\hat\lambda,\hat\nu)\|\)——Lagrangian 梯度的范数，最优时应该为零；
\item
  对偶间隙：\(f_0(\hat x)-g(\hat\lambda,\hat\nu)\)——原始值与对偶值之差，强对偶下应收敛到零。
\end{itemize}

这些残差必须结合问题的实际尺度来判读。\(10^{-6}\) 的约束违反对以微米为单位的精密加工是严格的，对以亿元为单位的财政约束连浮点误差级别都不到。变量缩放、约束归一化和目标函数的量级会直接影响残差的数值表现。更隐蔽的是：很多求解器在内部对问题做了自动缩放预处理，报告的残差可能是缩放后空间的残差而非原始尺度的残差。两个不同求解器都报告``KKT 容差满足''，背后比较的可能是完全不同的量——一个是缩放后的驻点残差，另一个是原始尺度的互补残差。

KKT 条件的最终角色是一份交接文件：它把``寻找最优点''这个优化问题，转化为``求解一组方程和不等式''这个代数问题。四组条件各自承担一部分几何约束——可行性划定允许区域、对偶性固定惩罚方向、互补松弛连接约束与价格的对应、驻点条件说合力为零。任何一条的缺失都意味着你还没有完全刻画最优性。
